% 1. Wprowadzenie
\section{Wprowadzenie}\label{sec:wprowadzenie}

\begin{instrukcja}[zalecane]
Krótko: czym jest aplikacja, jaki problem rozwiązuje i co zespół zrobił w ramach projektu. Wypiszcie,
czego świadomie nie robiliście (poza zakresem) - to pokazuje, że trzymaliście się minimum z karty
projektu. Pół strony wystarczy.
\end{instrukcja}

\subsection{Cel projektu}
W sesji student ma kilkanaście terminów z różnych przedmiotów rozrzuconych po wiadomościach,
kalendarzach i notatkach. Planer zbiera je w jednym miejscu, na starcie pokazuje najbliższy termin
z liczbą dni do niego, a na liście wyróżnia terminy pilne. Dzięki temu student szybko sprawdza,
co ma najbliżej i jaki materiał obowiązuje.

\subsection{Zakres projektu}
Zakres minimalny z karty projektu i jego realizacja:

\begin{table}[H]\centering
\caption{Realizacja zakresu minimalnego}\label{tab:zakres}
\begin{tabularx}{\linewidth}{@{}p{0.6cm}p{5.4cm}p{2.6cm}L@{}}
\naglowektabeli \nt{\#} & \nt{Element zakresu} & \nt{Stan} & \nt{Gdzie w aplikacji} \\
1 & ekran startowy & \tak & Start (\texttt{/}) \\
2 & lista elementów & \tak & Terminy (\texttt{/terminy}) \\
3 & ekran szczegółów & \tak & Szczegóły (\texttt{/termin?id=}) i panel obok listy \\
4 & formularz z walidacją & \czesciowo & Dodaj termin (\texttt{/dodaj}); reguły w \texttt{validate.ts}, formularz w budowie \\
5 & nawigacja łącząca co najmniej 4 ekrany & \tak & stałe menu: Start, Terminy, Dodaj termin; okruszki \\
6 & wersja na wąski ekran & \tak & dolny pasek nawigacji, jedna kolumna (rozdział \ref{sec:waski}) \\
\bottomrule
\end{tabularx}
\end{table}

\textbf{Poza zakresem} (świadomie pominięte): konta i logowanie, synchronizacja między urządzeniami, serwer, wysyłanie powiadomień (przypomnienie jest tylko zapisywane przy terminie), import z USOS.
